% Standalone English evaluation-design draft; no measured results are asserted.
% Compile: latexmk -xelatex WorldEdge-Table1-proposal.tex
% Research memo: docs/research/20261020-worldedge-table1/README.md
\documentclass[10pt,a4paper]{article}
\usepackage[margin=19mm]{geometry}
\usepackage[T1]{fontenc}
\usepackage{lmodern,amsmath,amssymb,booktabs,tabularx,array,microtype}
\usepackage{xcolor}
\usepackage[colorlinks=true,urlcolor=blue!50!black,linkcolor=blue!50!black]{hyperref}
\usepackage{placeins}
\setlength{\tabcolsep}{4pt}
\renewcommand{\arraystretch}{1.18}
\setlength{\emergencystretch}{2em}
\newcommand{\NE}{\textit{NE}}
\newcommand{\up}{\ensuremath{\uparrow}}
\newcommand{\head}[1]{\shortstack{#1}}
\title{WorldEdge: Proposed Main Evaluation Table}
\author{}
\date{}
\begin{document}
\maketitle
\vspace{-2em}
This design evaluates task-conditioned, persistent world generation and its utility for a fixed embodied executor. It separates world verification from actual task execution. All entries below are unmeasured placeholders; the candidate baselines require a documented capability and adapter audit before inclusion.

\begin{table}[htbp]
\centering\small
\caption{\textbf{World quality and downstream task execution: proposed design.} All methods receive the same preregistered task specifications, paired world seeds, robot embodiment, executor, observation interface, and budgets. Rates are percentages. \NE{} means not evaluated; it is not zero. The external-system and internal-ablation groups answer different comparisons and must retain separate interpretations.}
\label{tab:main}
\textbf{(a) World generation and independent verification}\par\smallskip
\begin{tabularx}{\linewidth}{@{}Xrrrrrr@{}}
\toprule
World generator & \head{Ready\up} & \head{Spec\up} & \head{Physics\up} & \head{Access\up} & \head{Interact\up} & \head{Persist\up} \\
\midrule
\multicolumn{7}{@{}l}{\textit{Candidate complete systems; common-interface evaluation}} \\
SceneSmith (adapted) & \NE & \NE & \NE & \NE & \NE & \NE \\
EmbodiedGen V2 (adapted) & \NE & \NE & \NE & \NE & \NE & \NE \\
SceneCode (adapted) & \NE & \NE & \NE & \NE & \NE & \NE \\
WorldEdge (full) & \NE & \NE & \NE & \NE & \NE & \NE \\
\midrule
\multicolumn{7}{@{}l}{\textit{WorldEdge ablations; shared assets and remaining components}} \\
One-shot generation & \NE & \NE & \NE & \NE & \NE & \NE \\
Without layout refinement & \NE & \NE & \NE & \NE & \NE & \NE \\
Without articulated assets & \NE & \NE & \NE & \NE & \NE & \NE \\
Without state persistence & \NE & \NE & \NE & \NE & \NE & \NE \\
\bottomrule
\end{tabularx}

\medskip
\textbf{(b) Downstream task success with a fixed agent}\par\smallskip
\begin{tabularx}{\linewidth}{@{}Xrrrrr@{}}
\toprule
World generator & \head{Navigation\\SR\up} & \head{Interaction\\SR\up} & \head{Cross-room\\SR\up} & \head{Revisit-state\\SR\up} & \head{Macro E2E\\SR [95\% CI]\up} \\
\midrule
\multicolumn{6}{@{}l}{\textit{Candidate complete systems; same executor across rows}} \\
SceneSmith (adapted) & \NE & \NE & \NE & \NE & \NE \\
EmbodiedGen V2 (adapted) & \NE & \NE & \NE & \NE & \NE \\
SceneCode (adapted) & \NE & \NE & \NE & \NE & \NE \\
WorldEdge (full) & \NE & \NE & \NE & \NE & \NE \\
\midrule
\multicolumn{6}{@{}l}{\textit{WorldEdge ablations; the task requirements remain unchanged}} \\
One-shot generation & \NE & \NE & \NE & \NE & \NE \\
Without layout refinement & \NE & \NE & \NE & \NE & \NE \\
Without articulated assets & \NE & \NE & \NE & \NE & \NE \\
Without state persistence & \NE & \NE & \NE & \NE & \NE \\
\bottomrule
\end{tabularx}
\par\smallskip
\begin{minipage}{\linewidth}\footnotesize
\textbf{Coverage.} Publish materialized/requested, admitted/requested, and executed/requested counts for every row in an accompanying coverage ledger. Access, interaction, and persistence use fixed, task-defined eligible sets, with identical denominators across methods.\par
\textbf{Scope.} A method not evaluated on a mandatory task family cannot receive a comparable Macro E2E score. A method evaluated in that family and failing to support a requirement receives a failure, not N/A. Reference/scripted executors and native-versus-adapted results must be labeled separately.
\end{minipage}
\end{table}
\FloatBarrier

\section{World metrics}
\textbf{Spec} is the fraction satisfying every required object identity, category/count, initial relation and state, protected constraint, and difficulty bound. Task goals must be unsatisfied initially. Generation must not simplify or replace the frozen specification.

\textbf{Physics} requires actual simulator loading and prescribed collision, support, boundary, and short-duration stability checks. Geometric thresholds, simulation duration, robot embodiment, and collision geometry are fixed on development data. Failed objects remain in the accounting; deleting unstable objects cannot improve the score without a specification penalty.

\textbf{Access} requires collision-aware routes from the designated start to every required observation or interaction region for the specified robot. It is a world-level prerequisite check, not navigation-policy success. A connected room graph or a large connected floor area alone is insufficient.

\textbf{Interact} requires a physical trace from fixed, independent probes of task-relevant affordances, such as reaching a prescribed articulation state or placing a specified object stably. Probes run on a reset or clone of the same delivered world. Passing primitive probes does not certify solvability of an entire long-horizon task; a failed probe leaves capability uncertified rather than proving impossibility.

\textbf{Persist} checks protected state after prescribed edits, expansion, departure, and re-entry or reload. Compare against the authoritative event history: preserve legal changes and reject unintended changes to object IDs, assets, poses, joint states, containment, and existing topology. Use a world-level all-required-constraints pass rate, with state and pose tolerances declared in advance.

\textbf{Ready} is the jointly verified yield, computed per request:
\begin{equation}
 \mathrm{Ready}=\frac{1}{N}\sum_{i=1}^{N}
 M_i\,V_i^{\mathrm{spec}}V_i^{\mathrm{phys}}V_i^{\mathrm{access}}
 V_i^{\mathrm{interact}}V_i^{\mathrm{persist}},
\end{equation}
where $M_i$ indicates materialization and each $V_i$ indicates a verified pass. Only a condition declared inapplicable by the task specification before generation is assigned one without a test. Missing measurements are not passes. Do not multiply the marginal column means. This is evidence of specified checks, not an oracle guarantee of complete task solvability.

\section{Agent tasks and success}
The task families are a proposed WorldEdge design. Navigation involves target search and arrival without manipulation. Interaction includes local pick/place and articulation, reported separately in the appendix. Cross-room tasks require transport or composition across rooms. Revisit-state tasks require changing a state, leaving or expanding the world, returning, and acting using the updated identity and state. For exclusive assignment, use the preregistered precedence: revisit-state, cross-room, local interaction, navigation. Action horizon and history lag are additional stratification variables.

Success must be verified from actual simulator trajectories, including identity, temporal dependencies, protected constraints, and final-state retention. A generated plan, a symbolic solver result, or a language-model success declaration is insufficient. Freeze the executor checkpoint, prompt, memory, sensors, controller, action abstractions, and budget across rows. Disclose any privileged observations, snap actions, or teleportation.

Panel (a) uses unique world-generation requests. Panel (b) assigns a fixed number $K$ of executor seeds to every generated world. For family $f$, $N_f$ counts all preregistered episode slots (world requests times $K$), and $S_f$ counts complete executions satisfying the frozen task contract. A generation failure contributes $K$ unsuccessful slots; these are not $K$ independent worlds. Report both world and episode counts. Then
\begin{equation}
 \mathrm{E2E\!\!-\!\!SR}_f=\frac{S_f}{N_f},\qquad
 \mathrm{Macro\ E2E\ SR}=\frac14\sum_{f=1}^{4}\frac{S_f}{N_f}.
\end{equation}
Generation, asset, initialization, specification, and genuine execution failures remain in the system denominator. Their causes are recorded separately and must not all be attributed to agent capability. With an actual mandatory admission gate admitting $A_f$ episode slots,
\begin{equation}
 \frac{S_f}{N_f}=\frac{A_f}{N_f}\frac{S_f}{A_f}.
\end{equation}
When $A_f=0$, E2E SR is zero and conditional SR is undefined (N/A). Conditional success is a diagnostic; it is not a substitute for the all-request endpoint. The identity holds within each family, not by multiplying macro-averaged terms. Independent certification is not necessarily a runtime gate, and the agent can succeed on an uncertified instance.

\section{Comparison and reporting protocol}
Pair methods on externally fixed task cards and world seeds. Prevent trivially easy or already-completed worlds through frozen initial conditions, affordances, density bounds, dependency depth, and history requirements. Audit realized difficulty without selecting a favorable subset after observing results. Primary intervals retain nested executor seeds and account for house/world and task-template dependence. When task templates and world families are crossed, use a design-appropriate multiway cluster procedure or hierarchical model rather than one-way episode resampling. Preregister one primary baseline contrast on Macro E2E SR; handle the multiplicity of any additional confirmatory claims.

Unrelated infrastructure failures receive a fixed retry procedure and an explicit unresolved count. If $U_f$ episode outcomes remain unknown in family $f$, report $[S_f/N_f,(S_f+U_f)/N_f]$. The Macro E2E identification bounds are $[\frac14\sum_f S_f/N_f,\frac14\sum_f(S_f+U_f)/N_f]$, separate from sampling confidence intervals. Do not silently drop unresolved outcomes or declare a completed comparison. Determine test-set size from an independent pilot and cluster-aware power or precision planning.

A separate learning-utility experiment should compare no adaptation, random curriculum, C1, C2, and C1+C2 under matched experience budgets and the same frozen held-out worlds. Table~\ref{tab:main} measures world-builder utility under a fixed policy; it does not establish transferable learning gains. Report visual quality, component diagnostics, partial goal completion, navigation SPL, latency, total cost, and growth curves separately.

\paragraph{Primary methodological references.}
\href{https://arxiv.org/html/2503.14756v3}{SceneEval} separates specification fidelity from spatial plausibility;
\href{https://proceedings.mlr.press/v306/pfaff26a.html}{SceneSmith} evaluates physics and scene quality;
\href{https://arxiv.org/html/2312.09067}{Holodeck} separates layout quality from downstream navigation;
\href{https://arxiv.org/html/2411.00081}{PARTNR} evaluates trajectory propositions and constraints;
\href{https://arxiv.org/pdf/1912.01734}{ALFRED} distinguishes task success from partial goal completion;
\href{https://arxiv.org/html/2505.16928}{Infinity-THOR} separates memory questions from online actions;
\href{https://arxiv.org/html/2012.02096}{PAIRED} motivates evaluating curriculum utility on held-out tasks. The proposed metric aggregation, task-family partition, and persistence probes are WorldEdge-specific design choices.

\paragraph{Current measurement boundary.}
The audited WorldEdge pipeline contains world-generation, session, and symbolic-solver components, but a complete, comparable physical task-execution study on its generated mesh/articulated worlds was not found. Existing FastFill strict admission, symbolic completion, and route matching must not populate Panel (b). No new experimental measurements are reported in this document.
\end{document}
